\question[4]
Die Hilfe zu zwei Methoden lautet:
\begin{lstlisting}[language={}]
sort(self, /, *, key=None, reverse=False)
    Sort the list in ascending order and return None.

shuffle(x)
    Shuffle list x in place, and return None.
\end{lstlisting}
\begin{parts}
\part Was erscheint auf der Konsole?
\begin{lstlisting}
a = [3, 1, 2]
b = a.sort()
print(a, b)
\end{lstlisting}
\part Was bedeutet \texttt{in place}?
\part Eine Liste \texttt{karten} soll gemischt werden. Warum ist \texttt{karten = random.shuffle(karten)} falsch?
\end{parts}
\begin{solutionbox}{2cm}
a) \texttt{[1, 2, 3] None} \\
b) Die Liste selbst wird verändert, es entsteht keine neue Liste. \\
c) \texttt{shuffle} gibt \texttt{None} zurück, \texttt{karten} wäre danach \texttt{None}. Richtig ist \texttt{random.shuffle(karten)}.
\end{solutionbox}
